%%% Pacotes utilizados %%%

% Codificação do arquivo
\usepackage[utf8]{inputenc}

% Indenta o primeiro parágrafo de cada seção.
\usepackage{indentfirst}
\usepackage{pdfpages}
\usepackage{subcaption}
\usepackage{float}
% Para adicionar link
\usepackage{hyperref}

% pacote cancel
\usepackage{cancel}

% Mapear caracteres especiais no PDF
\usepackage{cmap}

% Inserir PDF
\usepackage{pdfpages}

% Codificação da fonte

% ----------------------------------------------------------
\usepackage[T1]{fontenc}
\renewcommand{\rmdefault}{phv} % Arial
\renewcommand{\sfdefault}{phv} % Arial
% ----------------------------------------------------------

% Essencial para colocar funções e outros símbolos matemáticos
%\usepackage{amsmath,amssymb,amsfonts,textcomp}
\usepackage{amsmath}
\usepackage{mathtools}
\usepackage{amssymb,amsfonts,textcomp}
\usepackage{float}
\usepackage{calligra}
\usepackage{calrsfs}
\DeclareFontShape{T1}{calligra}{m}{n}{<->s*[1.4]callig15}{}
\DeclareMathAlphabet{\mathcalligra}{T1}{calligra}{m}{n}
\usepackage[mathscr]{euscript}
\usepackage{eqnarray}

% Para definir espaçamento entre as linhas
\usepackage{setspace}

% Espaçamento do texto para o frame
\setlength{\fboxsep}{1em} % 

%% Elementos Gráficos
% Para incluir figuras (pacote extendido)
\usepackage[]{graphicx}

%% Suporte a cores
\usepackage{color}
% Os argumentos declaram nomes novos, como Cyan e Crimson
% (ver documentação do pacote).
%\usepackage[usenames,dvipsnames,svgnames]{xcolor}

% Caso queira guardar as figuras em uma pasta separada
% (descomente e) defina o caminho para o diretório:
\graphicspath{{Imagens/}}
\usepackage{tabu}
\usepackage{bm}
\usepackage[labelfont=bf]{caption}

% Customizar as legendas de figuras e tabelas
\usepackage{caption}

% Customizar o tamanho de tabelas grandes
\usepackage{adjustbox}

%pacotes novos
\usepackage{siunitx,booktabs} 

% Criar ambientes com 2 ou mais linhas
\usepackage{multirow}

% Criar ambientes com 2 ou mais colunas
\usepackage{multicol}

%% Tabelas
% Elementos extras para formatação de tabelas
\usepackage{array}

% Tabelas com qualidade de publicação
\usepackage{booktabs}

% Para criar tabelas maiores que uma página
\usepackage{longtable}

% adicionar tabelas e figuras como landscape
\usepackage{lscape}

%% Lista de Abreviações
% Cria lista de abreviações
\usepackage[notintoc,portuguese]{nomencl}
\makenomenclature

%% Notas de rodapé
% Lidar com notas de rodapé em diversas situações
\usepackage{footnote}

% Notas criadas nas tabelas ficam no fim das tabelas
\makesavenoteenv{tabular}

% Conta o número de páginas
\usepackage{lastpage}

\usepackage{verbatim} %Comenta em bloco {02/01/2023}

% Verbatim com math mode
\usepackage{listings}
\lstset{
  basicstyle=\ttfamily,
  mathescape
}

% ----------------------------------------------------------
% Pacotes de citações
% ----------------------------------------------------------
\usepackage[brazilian,hyperpageref]{backref}	 % Paginas com as citações na bibl
\usepackage[alf]{abntex2cite}	% Citações padrão ABNT




%\bibliographystyle{apalike2345}
%\bibliographystyle{apalike}


% ----------------------------------------------------------
% CONFIGURAÇÕES DE PACOTES
% ----------------------------------------------------------

% ----------------------------------------------------------
% Configurações do pacote backref
% Usado sem a opção hyperpageref de backref
\renewcommand{\backrefpagesname}{Citado na(s) página(s):~}
% Texto padrão antes do número das páginas
\renewcommand{\backref}{}
% Define os textos da citação
\renewcommand*{\backrefalt}[4]{
	\ifcase #1 %
		Nenhuma citação no texto. % 
	\or
		Citado na página #2.%
	\else
		Citado #1 vezes nas páginas #2.%
	\fi}%
% ----------------------------------------------------------

%% Pontuação e unidades
% Posicionar inteligentemente a vírgula como separador decimal
\usepackage{icomma}

\usepackage{subfig} %Subfiguras

% Formatar as unidades com as distâncias corretas
% \usepackage[tight]{units} % 02/01/2023


% Dados do projeto
\newcommand{\nomedoaluno}{José Pereira Ramos Junior}

%% Comandos customizados

\newcommand{\mychapter}[2]{
    \setcounter{chapter}{#1}
    \setcounter{section}{0}
    \chapter*{#2}
    \addcontentsline{toc}{chapter}{#2}
}

% TIPOGRAFIA CUSTOMIZADA
% ----------------------------------------------------------
\renewcommand{\ABNTEXchapterfontsize}{\Large\bfseries}
\renewcommand{\ABNTEXsectionfontsize}{\Large\bfseries}
\renewcommand{\ABNTEXsubsectionfontsize}{\large}
\renewcommand{\ABNTEXsubsubsectionfontsize}{\large}
% ----------------------------------------------------------

% Espécie e abreviação
\newcommand{\subde}{\emph{Clypeaster subdepressus}}
\newcommand{\subsus}{\emph{C.~subdepressus}}

\titulo{Notas de Estudo de Viga}
\autor{\nomedoaluno}
\local{Brasil}
\data{2015}

% O tamanho do parágrafo é dado por:
\setlength{\parindent}{1.3cm}

% Controle do espaçamento entre um parágrafo e outro:
\setlength{\parskip}{0.2cm}  % tente também \onelineskip

% alterando o aspecto da cor azul
\definecolor{blue}{RGB}{41,5,195}

\makeatletter
\renewcommand{\@chapapp}{} % Not necessary...
\newenvironment{chapquote}[2][2em]
  {\setlength{\@tempdima}{#1}%
   \def\chapquote@author{#2}%
   \parshape 1 \@tempdima \dimexpr\textwidth-2\@tempdima\relax%
   \itshape}
  {\par\normalfont\hfill--\ \chapquote@author\hspace*{\@tempdima}\par\bigskip}

\hypersetup{
     	%pagebackref=true,
		pdftitle={\@title}, 
		pdfauthor={\@author},
    	pdfsubject={\imprimirpreambulo},
	    pdfcreator={LaTeX with abnTeX2},
		pdfkeywords={abnt}{latex}{abntex}{abntex2}{trabalho acadêmico}, 
		colorlinks=true,       		% false: boxed links; true: colored links
    	linkcolor=black,          	% color of internal links
    	citecolor=black,        	% color of links to bibliography
    	filecolor=black,      		% color of file links
		urlcolor=black,
		bookmarksdepth=4
}
\makeatother
% ----------------------------------------------------------
